% !TeX root = ../动手学强化学习笔记.tex
% 由主文档通过 \input 载入；请编译主文档。
\setcounter{chapter}{6} % 保留原书章号。
\chapter{蒙特卡洛方法}
\label{ch:rl-mc}

模型未知时，不能直接计算贝尔曼方程中的全部后继期望。蒙特卡洛方法（Monte Carlo method，MC）用实际经历的完整回报估计价值：它不需要转移表，也不依赖对后继价值的当前猜测，但必须等待足够完整的结果。

\section{回合采样与蒙特卡洛预测}

考虑固定策略下独立重置的回合，终止时刻为 $T$。从 $G_T=0$ 开始，逆序计算 $G_t=R_{t+1}+\gamma G_{t+1}$。首次访问（first-visit）MC 每回合只使用某状态第一次出现时的回报；每次访问（every-visit）MC 使用该状态所有出现时的回报。两者都可按式\eqref{eq:rl-incremental-mean}更新均值，但使用的样本集合不同。

\begin{example}[首次访问与每次访问]
某回合依次产生 $S_0=s,R_1=1,S_1=s,R_2=2,S_2=\text{终止}$，取 $\gamma=1$。两个回报为 $G_0=3,G_1=2$。仅有这一回合时，首次访问估计为 $3$，每次访问估计为 $(3+2)/2=2.5$；它们不是对同一组数字执行不同平均公式，而是选取了不同样本。
\end{example}

\begin{proposition}[首次访问 MC 的一致性]
固定策略和独立同分布重置规则下，若状态 $s$ 在一个回合中以正概率被访问，首次访问后的回报可积，则在无限多个回合中，首次访问 MC 的估计几乎必然收敛到 $v_\pi(s)$。
\end{proposition}
\begin{proof}
含有访问的回合以正概率重复出现，几乎必然有无限多个。由马尔可夫性质，在每个回合第一次到达 $s$ 后，后续回报的条件分布都是从 $s$ 出发执行 $\pi$ 的回报分布。独立回合给出独立同分布的这些回报，应用强大数定律即得结论。
\end{proof}

每次访问样本在回合内相关。其一致性可通过按回合汇总回报和访问次数的比值分析，在有限期望访问次数及相应可积条件下成立；不能简单声称同一回合中的全部样本独立。

\section{动作价值估计与蒙特卡洛控制}

已知 $v_\pi(s)$ 但不知道环境转移时，不能直接计算每个动作的后果，因此控制通常估计 $q_\pi(s,a)$。把首次访问对象由状态换成状态动作对，即可得到动作价值 MC。\emph{探索起始}（exploring starts）假定每个需要学习的状态动作对都有正概率作为起点；这在模拟模型中有时可行，在现实系统中未必可行。

实用的软策略控制保持动作探索。初始化 $Q$、各对访问次数和一个 $\varepsilon$-贪心策略；每回合按当前策略采样，逆序计算回报，对本回合首次出现的状态动作对更新均值，再令策略对新 $Q$ 作 $\varepsilon$-贪心选择。终止条件通常是交互预算耗尽，输出价值估计与指定评价策略。固定 $\varepsilon$ 保持软行为，减小 $\varepsilon$ 则需要确保状态动作对仍被充分访问；仅把 $\varepsilon_t$ 写成递减序列不自动保证整个 MDP 中的覆盖。

控制过程的策略不断改变，历史回报来自不同策略，因而固定策略一致性命题不能直接当作控制收敛证明。理想化分析可采用每轮获得精确或充分准确评估后再改进；一般逐回合控制的保证必须针对具体探索与更新条件陈述。

\section{偏差、方差与回合长度}

在固定策略与可积回报下，完整回报 $G_t$ 是从给定状态预测 $v_\pi$ 的无偏样本，但其方差可能很大。长回合累积多次随机事件，稀疏终点奖励又使大量失败轨迹提供相同数值。MC 不用其他价值估计自举，代价是无法在回合完成前利用尚未观测的最终回报。

若 $|R_t|\leq R_{\max}$、$\gamma<1$，只保留未来 $h$ 步产生的尾项绝对误差至多为 $\gamma^hR_{\max}/(1-\gamma)$。这个界解释截断误差，并不使截断回报变成无偏完整回报。若任务真实终止，尾项确为零；若只是采样停止，则应承认截断偏差或引入后继价值估计，后者通向 TD 方法。

\section{习题与实验}
\begin{enumerate}
\item 对回合 $s\to u\to s\to\text{终止}$，奖励依次为 $0,1,2$，取 $\gamma=1/2$，计算全部回报及 $s$ 的两种 MC 估计。
\item 随机游走状态为 $0,1,2,3,4$，每步等概率左右移动，$0,4$ 为终止；到达 $4$ 奖励 $1$，其余奖励为零。无折扣时证明 $v(i)=i/4$，并用独立回合 MC 验证。
\item 比较固定回合数与固定交互步数下的算法评价。为什么不同平均回合长度会使两种预算不等价？
\end{enumerate}





